\documentclass[11pt]{article}

\usepackage[margin=1in]{geometry}
\usepackage[T1]{fontenc}
\usepackage{lmodern}
\usepackage{booktabs}
\usepackage{tabularx}
\usepackage{enumitem}
\usepackage{url}
\usepackage{hyperref}

\newcommand{\todo}[1]{\textbf{[TODO: #1]}}
\newcommand{\code}[1]{\texttt{#1}}

\title{LLM-Generated Binary Proofs: \\
       OSS-Verify on LPS and glibc on Picinae}
\author{Charles Averill \and Ilan Buzzetti \and Shreya Soman \and Kevin Hamlen \\[1ex]
        The University of Texas at Dallas}
\date{Draft, \today}

\begin{document}
\sloppy
\maketitle

\begin{abstract}
In October 2026, Anthropic and Theorem reported that language-model agents
had written Rocq proofs about compiled RISC-V binaries of OSS-Fuzz targets.
The proofs show that sanitizer builds of those targets never fault and behave
like their plain builds~\cite{anthropic-post,oss-verify}.  This note compares
that work with our own experiment, in which Claude Opus~5.5, running in
Claude Code, wrote Picinae proofs about functions lifted from the x86-64
glibc~2.35 shared library.  The two efforts share a proof assistant and a
binary-level starting point.  They differ in almost everything else: the
semantic model and how the program enters it, the kind of property proved,
how specifications are obtained, the scale, and what has to be trusted.  We
also discuss the lifting workflow that a collaborator has written up as an
agent skill, and how it relates to the lifter changes made in our
experiment.
\end{abstract}

\section{Introduction}

The Anthropic post~\cite{anthropic-post} argues that language models make
formal verification of real software feasible at a much larger scale than
before.  Its headline result is 239 fuzz-test entry points in 123 open-source
libraries, for which Claude proved that ``there are no crashes on any input
(up to 64\,MB)''.  The proofs are stated over compiled RISC-V instructions
under a formal model of a Linux process.  The public repository from Theorem,
\code{theorem-labs/oss-verify-proofs}~\cite{oss-verify}, holds a checked
release of part of that work.

Our project started from that post.  We wanted to know how well the same kind
of agent does on a different binary-verification platform, Picinae, with a
different kind of specification: per-function functional correctness taken
from the man pages, not crash-freedom of whole programs.  The work so far is
in \code{Examples/claude} of the Picinae repository, on the branch
\code{llm\_generated\_proofs}.

Section~\ref{sec:glance} summarizes the two efforts.
Section~\ref{sec:models} compares the semantic models,
Section~\ref{sec:what} what the theorems say,
Section~\ref{sec:hard} what makes our proofs hard, and
Section~\ref{sec:method} the methodology.
Section~\ref{sec:lifting} discusses lifting, including the collaborator's
skill.  Section~\ref{sec:discussion} draws lessons.

\section{The two efforts at a glance}\label{sec:glance}

\begin{table}[h]
\centering\small
\begin{tabularx}{\textwidth}{@{}lXX@{}}
\toprule
 & OSS-Verify (Anthropic / Theorem) & This project (Picinae) \\
\midrule
Proof assistant & Rocq 9.2.0 & Rocq 9.1 \\
Target code & whole static RISC-V (RV64GC) executables: OSS-Fuzz harness, library, musl, sanitizer runtime; vector code paths avoided & individual functions of the x86-64 glibc 2.35 shared object, including the hand-written AVX2 assembly the machine actually runs, plus direct callees where needed \\
Program representation & ELF bytes loaded into a memory image, decoded by the Rocq ISA semantics & Picinae IL produced by the Ghidra p-code lifter \texttt{pcode2coq} \\
Semantics & LPS: RISC-V instructions plus Linux process semantics (system calls, start-up, budgets) & Picinae amd64 IL semantics; no OS model \\
Property & sanitizer build never errs and agrees with the plain build (also halting, bound and decidability shapes) & functional correctness against the man page, plus frame and register preservation \\
Correctness kind & safety over finite traces (agreement); total for the halts shape & partial correctness (holds if the function returns) \\
Specifications & generated mechanically from fuzz harnesses, from a fixed set of statement shapes & written by the model per function, from the man page and glibc behavior \\
Scale & post: 239 entry points, 123 libraries, about $1.8\times10^8$ lines of Rocq; repository release: 138 statements proved for 60 entry points in 25 libraries & 78 proof files covering about 80 glibc entry points (among them the AVX2 \code{memmove}/\code{memcpy}, \code{strlen}, \code{strchr} and the short-string case of \code{strcmp}), about $1.9\times10^4$ lines of proofs and libraries \\
Agents & many autonomous agents with a shared workspace, over a few weeks & a single Claude Code session, with occasional human direction, over about three days \\
Unfinished proofs & not allowed: \code{Admitted} refused, axioms audited & four main theorems left \code{Admitted} because \code{Qed} is too slow, each with a note \\
\bottomrule
\end{tabularx}
\caption{Summary.  The OSS-Verify figures come from the post~\cite{anthropic-post} and the repository README~\cite{oss-verify}; the two do not cover the same set of proofs (Section~\ref{sec:scale}).}
\end{table}

\section{Semantic models}\label{sec:models}

\subsection{How the program enters the logic}

In OSS-Verify the program is its bytes.  The tool \code{elf-snapshot} turns
an ELF file into a memory image (loadable segments, entry point, program
headers).  An OCaml plugin, \code{rocq-read-file}, reads the image into Rocq
at compile time, and the RISC-V semantics in LPS decodes and executes
instructions from that image.  No disassembler or lifter is involved.  The
trusted path from binary to logic is the ELF loader and the image reader,
and \code{Images.v} checks by computation that each image is well formed.

In Picinae the program is a Rocq function from addresses to IL statements.
It is produced outside the logic by Ghidra (disassembly plus the SLEIGH
specification of each instruction as p-code) and by \code{pcode2coq}, which
translates p-code to Picinae IL.  So the disassembler, the SLEIGH
specifications and the translator are in the trusted base.  The experiment
underlined this risk: several translator bugs gave unsound lifts until we
fixed them.  For example, partial-register writes clobbered the rest of the
register, \code{LZCOUNT} was lifted as width minus popcount, and flag
booleans had the wrong width (Section~\ref{sec:lifting}).  A Picinae proof
is a proof about the lifted IL; it is about the binary only to the extent
that the lift is right.

Lifting does have advantages.  The IL is architecture-neutral, so one
symbolic interpreter and one simplifier (\code{psimpl}) serve every ISA
Picinae supports.  Proofs read in terms of IL variables and expressions, not
instruction encodings.  And because lifting works one function at a time,
small proofs stay small.

\subsection{Instruction semantics}

LPS implements RV64GC from the RISC-V specification, as an executable
semantics inside Rocq.  The post adds that for the zlib case study Claude
also wrote a formal x86 ISA semantics.  \todo{Check whether that x86
semantics is published, and how its fidelity was checked.}  Picinae's amd64
semantics is SLEIGH's, as translated.  Some instructions are modeled only
approximately:
\begin{itemize}[nosep]
\item the \code{LOCK} prefix lifts as a no-op, so atomicity is assumed, not
      modeled;
\item bits 256--511 of the ZMM registers are not modeled; AVX/AVX2 code
      only ever zeroes them;
\item SLEIGH user-defined operations that Ghidra leaves opaque (for example
      \code{vpcmpeqb\_avx2} and \code{vpminub\_avx2}) had to be spelled out
      byte by byte in the lifter.
\end{itemize}

\subsection{Memory, faults and the environment}

This is the most important semantic difference.

LPS models a Linux process.  Memory is mapped by the loader and by
\code{mmap}/\code{brk}, and an access outside it is an \emph{error}.  An
instruction or system call that LPS does not interpret is also an error
(\code{unmodeled}), and so are \code{abort()} (failed \code{assert}s and
sanitizer reports) and illegal instructions (failed UBSan checks).  The
statements forbid every error, so a proof rules them all out.  The
environment is nondeterministic, in the way a sound model of Linux must be:
\code{read} may return any nonempty prefix, \code{mmap} any free aligned
address, and some system calls may fail spuriously.  Proofs must cover every
such choice.  The model is single-threaded, runs no signal handlers, and has
a simplified file system.

Picinae's memory is a total function from addresses to bytes.  A load or
store has a premise that the address is readable or writable (bits of the
pseudo-variables \code{A\_READ}/\code{A\_WRITE}).  When the premise fails,
the access has no successor state; it is not an error.  Picinae theorems are
partial-correctness statements over the traces that exist, so a run that
would fault simply has no continuation, and the theorem says nothing about
it.  As a result, our theorems do \emph{not} establish memory safety.  They
say that \emph{if} the function runs to its return, the result and the final
memory are as specified.  Picinae has no OS model; our functions make no
system calls (the mutex functions reach the futex system call only inside a
callee, which they assume through a contract).

Picinae's memory model also lets the vectorized string functions read past
the end of a string without any special argument.  For example,
\code{\_\_strlen\_avx2} reads whole 32-byte and 128-byte blocks.  In LPS the
same reads must be shown to stay inside mapped pages, which is exactly the
page-cross reasoning the glibc code does.  Picinae proofs can skip that
reasoning, but they do not get its benefit either.

\subsection{Concurrency}

Both models are single-threaded.  For us this matters most for
\code{pthread\_mutex\_lock} and \code{pthread\_mutex\_unlock}, which rank
first and third among libc functions in our Firefox profile.  We prove their
sequential behavior for every mutex kind, assuming \code{LOCK}-prefixed
instructions are atomic and no other thread runs during the call.  That is
a real limitation for functions whose whole purpose is concurrency.

\section{What gets proved}\label{sec:what}

\subsection{OSS-Verify statement shapes}

Every OSS-Verify statement applies one of a few \emph{shapes} (in
\code{theories/Shapes.v}) to the memory images of a harness's builds.  Each
shape is proved, by \code{reflexivity}, equal to an LPS definition:
\begin{itemize}[nosep]
\item \textbf{halts:} every execution of a build exits with an allowed
      status, under a memory-budget premise;
\item \textbf{agrees:} on every input where the plain build stays within its
      budget, neither the plain build nor the sanitizer build errs, and the
      two have the same set of finite traces (stream bytes, exit status,
      final file contents);
\item \textbf{disagrees at most $b$:} under a stated input distribution, the
      builds disagree with probability at most $b$;
\item \textbf{agreement decidable:} some RISC-V program classifies exactly
      the agreeing inputs.
\end{itemize}
The published proofs are agreement statements.  According to the README, no
halts statement is proved in this release, so the proofs do not rule out
non-termination.  An agreement proof for an ASan build rules out every error
the sanitizer detects on every covered input, which is what ``no crashes''
means in the post.  It says nothing about functional correctness: a parser
that returns a wrong result without faulting satisfies its statement.

\subsection{Picinae glibc theorems}

Our theorems are functional specifications of individual functions.  Every
file starts with a header that states the function's semantics, quotes the
man page, names the model, and gives an English input/output specification.
Typical main properties:
\begin{itemize}[nosep]
\item exact results, for example \code{ffs} returns the 1-based index of the
      lowest set bit, and \code{rand\_r} equals glibc's C reference;
\item exact memory effects with a frame condition, for example memset sets
      \code{dest[0,n)} to \code{c} in every size class (including the
      \code{REP STOSB} path) and leaves every other byte unchanged;
\item for non-leaf functions, correctness relative to a contract for each
      callee: \code{bsearch}, \code{lfind} and \code{lsearch} through a
      function-pointer comparator, \code{strtok\_r}/\code{strsep} through
      PLT calls, \code{strdup} through an abstract heap for \code{malloc},
      and the mutex functions through \code{\_\_lll\_lock\_wait} and
      \code{\_\_lll\_lock\_wake}.
\end{itemize}
Each function also gets generic theorems: the lifted program is well typed,
code is not writable, callee-saved registers are preserved, and memory is
preserved where the function does not write it.

Two theorems record places where glibc does not meet its man page.
\code{sigisemptyset} tests only the low 32 bits of the set, so a set
containing only real-time signals is reported empty.  \code{a64l} does not
sign-extend its 32-bit result, although the man page says it does.  For each
we proved the binary's behavior and a concrete counterexample to the man
page, and confirmed the counterexample by running the system libc.
\todo{Decide whether to report these upstream; check whether they are
already known or documented.}

\subsection{Vacuity}

A theorem can be true for an unintended reason.  OSS-Verify guards against
this in a few ways.  Statements are generated from facts and rendered by a
tool anyone can rerun.  Each shape is shown equal to the LPS definition it
stands for.  \code{Images.v} checks that the images are well formed.
Statements are fingerprinted, so a proof is always of the published
statement.

Our main vacuity risk is in callee contracts.  A contract that no code can
satisfy makes any theorem that assumes it trivially true.  This happened:
an early \code{bsearch}/\code{lfind} contract omitted a conjunct of the
callee's entry invariant (the return address on the stack), and a
single-state trace refuted it.  Every contract now has a \code{*\_satisfiable}
theorem that exhibits code meeting it.  For some contracts the only witness
we found is code that never returns, which is weak evidence.

\section{What makes these proofs hard}\label{sec:hard}

OSS-Verify's achievement is breadth: hundreds of entry points, each a whole
program, checked mechanically.  Our proofs are few and small by comparison,
but each one has to get through difficulties that the OSS-Verify statements
do not ask for, or that its setup avoids.  This section lists them, with
the function where each one showed up.  We also note where OSS-Verify's
obligations are harder than ours.

\subsection{Correct results, byte by byte}

An agreement statement is satisfied by any program that neither faults nor
behaves differently under a sanitizer.  A \code{memcpy} that copies the wrong
bytes satisfies it.  Our \code{memcpy} theorem quantifies over every byte of
memory: \code{dest[0,n)} must equal \code{src[0,n)} afterwards and every
other byte must be unchanged.  glibc meets this with stores that overlap on
purpose:
\begin{itemize}[nosep]
\item two possibly overlapping head and tail moves for small sizes;
\item a 64-byte head, aligned 128-byte iterations and a saved tail for large
      sizes;
\item REP MOVSB starting at a 64-byte boundary, followed by stores of a saved
      head.
\end{itemize}
Every path needs a coverage argument that the union of the stored offset
ranges is exactly $[0,n)$.  For \code{memmove}, the stores must also be
shown not to clobber source bytes that have not yet been read.  The
OSS-Verify proofs do contain functional lemmas of this kind, for example a
specification of musl's \code{memcpy}, because whole-program reasoning needs
them.  But they are internal lemmas about compiled C, not theorems about the
code a production x86 system runs.

\subsection{Hand-optimized code with many strategies}

The functions Firefox calls most are glibc's hand-written AVX2 assembly,
selected at load time by IFUNC.  \code{\_\_memmove\_avx\_unaligned\_erms},
which serves both \code{memmove} and \code{memcpy}, has 351 instructions and
chooses among at least seven strategies by size, overlap, 4\,KiB aliasing
and CPU-feature flags.  One theorem over all of them was beyond what the
proof engine could hold, so we split it into three theorems by size range:
up to 256 bytes, the 128-byte loops, and REP MOVSB.  Each theorem refutes
the other paths' entries with \code{False} invariants.  The non-temporal path's dispatch and entry are proven, but its
interleaved multi-page store loops are admitted with a note; the
overlapping path above 256 bytes is still open.

\subsection{SIMD}

The string functions compare 32 bytes at a time (VPCMPEQB), combine vectors
with VPMINUB, collapse a comparison into a bit mask (VPMOVMSKB), and find the
first set bit (TZCNT, which the lifter turns into a 33-iteration loop).
\begin{itemize}
\item \textbf{Lifting.}  Ghidra leaves several AVX2 operations opaque.  The
      lifter had to spell them out lane by lane, and to merge per-lane
      register writes whose naive translation doubled the term size with every
      lane.
\item \textbf{Term size.}  Stepping one VPMOVMSKB symbolically executes
      about 200 IL moves and produces a 32-way disjunction over the bytes of
      memory.  The brute-force proof of a single mask fact took 106\,s.
      We instead wrote instruction-level lemmas: a Coq function generates
      each lifted instruction, the lemma about it is proved once by induction
      over the lanes, and a proof steps the instruction with a single lemma
      application.
\item \textbf{OSS-Verify avoids vector code.}  Its statements include a
      variant for binaries that take vector code paths when the processor
      advertises the RISC-V vector extension: the process is told the
      machine has no vector unit, so the scalar paths are what gets verified.
\end{itemize}

\subsection{Bit-level reasoning about positions}

The correctness of the string functions rests on index arithmetic over
masks.  For \code{strlen}:
\begin{itemize}[nosep]
\item bit $j$ of the mask is set iff byte $j$ of the block is zero;
\item TZCNT gives the first such $j$, and with the bytes already known to be
      nonzero that determines the length;
\item near a page end the mask of the aligned block is shifted right by
      $s \bmod 32$ with an \emph{arithmetic} shift (SARX), and a separate
      argument shows that the sign-replicated high bits never produce the
      first set bit;
\item the loop exit combines two 32-bit masks into one 64-bit value and uses
      a 64-bit TZCNT, which has to be split back into the two cases.
\end{itemize}
For \code{strchr}, a ``hit'' is a byte equal to 0 or to $c$; the loop
detects it as a zero byte of $\min(x, x \oplus c)$.  For \code{strcmp}, a
stop is a byte that differs between the two strings or terminates the first,
which the code computes as a zero lane of $\min(\mathrm{eq}(a,b), a)$.

\subsection{Configuration in memory}

The memcpy strategy depends on tunables that glibc sets at startup:
\begin{itemize}[nosep]
\item the REP MOVSB threshold and its stop threshold;
\item the non-temporal threshold;
\item a CPU-feature word that disables REP MOVSB for short distances.
\end{itemize}
The proof reads them as symbolic memory, so they become preconditions and
case splits.  Some preconditions encode facts glibc guarantees elsewhere,
for example that the REP MOVSB threshold is never below 256 bytes on AVX.

\subsection{x86 semantics}

Partial-register writes, the flags and their many encodings, REP prefixes
(lifted as instructions that jump to themselves, so REP STOSB and REP MOVSB
need their own invariants and the direction flag), and the LOCK prefix all
had to be handled.  Several needed lifter fixes for soundness
(Section~\ref{sec:lifting}).  Branches on signed comparisons force the flags
into the proof: \code{strcmp}'s page-safety test is a \code{JG} after a
\code{CMP}, so the proof must recover the signed relation from the lifted
\code{OF}, \code{SF} and \code{ZF} bit-expressions (the \code{OF} term turns
out to be identically zero), and it cannot abstract those flags away the way
the loop-only functions do.

\subsection{The proof engine at the limit}

A single glibc function can produce symbolic states larger than an
OSS-Verify proof ever needs in one goal, because its paths rejoin and its
values are wide.  The resulting failures were mostly engineering, not
mathematics:
\begin{itemize}[nosep]
\item tactics that compare memory terms (\code{assumption}, \code{rewrite}
      with unification variables, \code{discriminate}) unfold them and run
      out of memory;
\item \code{lia} runs out of memory on any hypothesis that mentions a load,
      and becomes very slow with several equations modulo $2^{64}$;
\item the kernel check at \code{Qed} can take far longer than the tactics.
\end{itemize}
Each needed a work-around, recorded in the README: syntactic tactics,
cleared contexts, explicit lemma instances, and extra cut points.

\subsection{Where OSS-Verify's obligations are harder}

The comparison cuts both ways.  OSS-Verify proofs must:
\begin{itemize}[nosep]
\item reason about whole programs, including start-up code, the C library
      and the sanitizer runtime;
\item handle a nondeterministic operating system: short reads, arbitrary
      \code{mmap} addresses, spurious failures;
\item show that every memory access stays inside mapped memory.
\end{itemize}
Our theorems are partial-correctness statements about single functions, and
Picinae's memory model lets them ignore faults.  The vectorized string
functions read past the end of the string, and in Picinae that needs no
argument.  The same reads in LPS would need the page-boundary reasoning the
glibc code is designed around.

\section{Methodology}\label{sec:method}

\subsection{Where specifications come from}

OSS-Verify turns testing infrastructure into theorems.  An OSS-Fuzz harness
already defines an entry point and an input format.  The statement says that
the harness behaves on all inputs the way fuzzing hopes it behaves on the
inputs it tried.  Producing a statement takes no understanding of the
library; this is what makes 20{,}164 statements over 2{,}148 harnesses
possible.  It also limits what can be said, because there is no notion of
the right output.

Our specifications are written by the model, function by function, from the
man page and from reading the binary.  They say much more, but they are a
second artifact that must be reviewed against the documentation.  For
example, the headers note where glibc departs from the man page
(\code{sigfillset} leaves out signals 32 and 33, which NPTL reserves), and
where a specification needed a precondition the man page does not mention
(\code{memcpy}'s loop theorem assumes \code{64 <= d}).  We have not yet
reviewed every specification against the C sources.
\todo{Add a specification-review pass, ideally by someone other than the
generating model.}

\subsection{Proof style}\label{sec:proofstyle}

The OSS-Verify proofs work at whole-program scale.  A typical proof file
tree includes interaction-tree models of the program, simulation arguments
between builds, and per-function lemmas for the C library.  For example,
the zstd \code{huf\_decompress} proof contains a specification of musl's
\code{memcpy} at its RISC-V entry address, with a non-overlap precondition
much like ours.  The proof scripts in the release are stored in a compressed
replay encoding (\code{KReplay}) with comments removed, so they are checkable
but not meant to be read.  The post reports a typical proof of about
600{,}000 lines for binaries with a median of 18{,}000 instructions.  The
repository's \code{proofs/} directory holds about $3.9\times10^7$ lines of
Rocq.

Picinae proofs follow the platform's Floyd--Hoare style.  The prover places
invariants at cut points (loop heads, joins, call sites).  The symbolic
interpreter (\code{x64\_step}) and the simplifier (\code{psimpl}) step
between them, and the remaining arithmetic and bit-vector goals are closed
by hand.  Proofs are written to be read, and the README records what we
learned about keeping them fast.  Much of the experiment's effort went into
three performance problems:
\begin{itemize}[nosep]
\item \textbf{Memory.}  Conversion near symbolic memory terms made plain
      \code{reflexivity}, \code{assumption} and \code{rewrite} blow up past
      8\,GB.  We replaced them with syntactic (\code{constr\_eq}-guarded)
      versions and added cut points at every join.
\item \textbf{Kernel time at \code{Qed}.}  Several tactic proofs finish in
      seconds or minutes, but their \code{Qed} runs for more than 30
      minutes.  The cause we found is the kernel unfolding \code{N.ltb}
      against power-of-two literals.  A \code{Strategy} declaration fixed some
      cases.  Four theorems remain \code{Admitted} with a note
      (\code{pthread\_mutex\_lock}, \code{\_\_memset\_avx2\_unaligned\_erms},
      \code{isinfl}, \code{rand\_r}).
\item \textbf{SIMD.}  Stepping an AVX2 lane instruction symbolically builds
      a 32-lane term; \code{VPMOVMSKB} alone executes about 200 IL moves.
      Our response was a library of instruction-level lemmas
      (\code{amd64\_avx2.v}).  Each lemma describes a whole lifted
      instruction through a Coq generator convertible to the lifter's
      output, and is proved once by induction over the lanes.
\end{itemize}
The OSS-Verify release does not allow unfinished work: \code{Admitted} and
new axioms are refused, and \code{Print Assumptions} is audited against an
allowlist.  We have no such audit yet; our libraries use
\code{FunctionalExtensionality}, an axiom that is also on their allowlist.
\todo{Add a \code{Print Assumptions} audit to the Makefile.}

\subsection{Agents and human involvement}

The post describes many autonomous agents with little human involvement,
sharing a workspace of finished proofs.  In a 36-task benchmark, agents that
shared proofs used 17\% fewer tokens, and some succeeded after failing on
their own.  When an agent found a bug, it patched the code, regenerated the
statement, and reproved it.  Patches were logged and audited, and the reward
was shaped to discourage changing the program to make a proof go through.
Proofs reportedly cost up to hundreds of millions of tokens.

Our experiment used one agent session (Claude Code with Claude Opus~5.5)
over about three days, with a human giving occasional direction.  Examples:
order the work from easiest to hardest; commit only when asked; leave a
proof \code{Admitted} with a note when its tactics are fast but its
\code{Qed} is slow; prioritize the functions that a profile of Firefox
visiting a website shows libc spending the most time in.  We did not patch
glibc.  Where the binary disagrees with its documentation, we prove what
the binary does.  \todo{Report token usage and wall-clock time per
function.}

\subsection{Scale}\label{sec:scale}

The two OSS-Verify sources give different numbers.  The post reports 239
entry points in 123 libraries.  The repository release reports 138 proved
statements, for 60 entry points in 25 libraries, out of 20{,}164 statements.
It also explains that proofs were ported to the release's binaries and LPS
version, which suggests the release is the subset ported so far.
\todo{Confirm with the OSS-Verify authors.}  Either way, the gap with us is
two to four orders of magnitude in proof volume.  A single OSS-Verify proof
is far larger than our whole project, mainly because each one reasons about
a whole static executable, including the C library and the sanitizer
runtime.

\section{Lifting}\label{sec:lifting}

\subsection{Our pipeline}

We lift from a Ghidra~11.4.2 project in which \code{libc.so.6} has been
imported and analyzed.  A headless driver,
\code{lifter/pcode2coq\_headless.py}, calls \code{pcode2coq}'s library entry
point with these settings:
\begin{itemize}[nosep]
\item relative addressing (entry point at 0) for single functions;
\item global addressing for a group of functions lifted into one program
      (\code{f1+f2=name}), so that direct calls between them resolve;
\item \code{strictFailure}, so unmodeled registers and operations abort the
      lift instead of becoming \code{Nop};
\item \code{typeCheck}, which appends a well-typedness theorem;
\item static-memory hypotheses switched off;
\item \code{canonicalTemps}, which numbers temporaries per instruction so
      that identical instructions lift to identical statements.
\end{itemize}
The AVX2 variants that the system actually runs are lifted by address.  The
lifter itself needed a series of soundness and capability fixes, kept on a
branch of \code{pcode2coq}:
\begin{itemize}[nosep]
\item partial-register writes merged into the base register;
\item exact varnode widths and correct flag widths;
\item \code{BOOL\_NEGATE} as \code{xor 1};
\item \code{LZCOUNT} lifted correctly;
\item intra-instruction branches, with the bit-scan loops of
      \code{BSF}/\code{TZCNT} bounded by \code{Rep};
\item AVX/AVX2 user operations spelled out;
\item per-lane vector writes merged into one move;
\item \code{LOCK} as a no-op.
\end{itemize}
\todo{Merge these fixes upstream; several are soundness bugs that affect
any amd64 lift.}

\subsection{The collaborator's lifting skill}

A collaborator working in parallel has written the lifting workflow as an
agent skill (\code{SKILL.md}).  It tells an agent to:
\begin{itemize}[nosep]
\item run \code{pcode2coq\_headless.py} from
      \code{Picinae/tools/ghidra-lifter} through Ghidra's \code{pyghidraRun}
      in headless mode, against a pre-existing Ghidra project;
\item select functions with a Python regular expression
      (\code{--functionRegex});
\item always pass \code{--globalAddressing}, ``so they can compose and be
      used together to reason about caller-callee relationships'';
\item pass \code{--typeCheck};
\item write the output to a subdirectory of \code{Examples} named after the
      binary, with one file per lifted function named after the function;
\item use \code{-noanalysis} only when the binary has already been analyzed;
\item leave \code{--prependProgName} and \code{--interleave} unset unless
      asked.
\end{itemize}

Compared with our pipeline:
\begin{itemize}
\item \textbf{The idea is the same.}  Both drive \code{pcode2coq}
      headlessly through PyGhidra, against a pre-analyzed project, with type
      checking on.  Writing the procedure down as a skill makes lifting
      something any agent can do reliably.  Our session had to rediscover
      the procedure (Ghidra and JDK versions, PyGhidra installation) before
      lifting anything.
\item \textbf{Addressing.}  The skill always uses global addressing.  We use
      relative addressing for leaf functions and global addressing only for
      groups that call each other.  Global addressing composes better:
      caller and callee proofs can share addresses, which is what
      \code{perform\_call}-style reasoning needs, and a function lifted
      alone can later join a group without being re-lifted.  Relative
      addressing keeps leaf proofs independent of where the library is
      loaded.  Given the skill's goal, global addressing everywhere seems
      the better default, and we should adopt it.
\item \textbf{Selection and layout.}  \code{--functionRegex} lifts many
      functions in one Ghidra start-up; our driver takes explicit names or
      addresses.  The skill's layout
      (\code{Examples/<binary>/<function>.v}) differs from ours
      (\code{Examples/claude/<function>\_amd64.v}).  We should agree on one
      before merging.
\item \textbf{Soundness options.}  The skill does not mention
      \code{strictFailure}.  Without it, the lifter emits \code{Nop} for
      unmodeled registers and operations, which silently changes the
      program; we recommend making it the default.  The skill also predates
      the static-memory, canonical-temporaries and AVX changes; whether its
      script includes them depends on which \code{pcode2coq} it uses.
      \todo{Check the collaborator's \code{tools/ghidra-lifter} version
      against our \code{pcode2coq} branch.}
\item \textbf{Variant selection.}  glibc picks string and memory routines
      at load time through IFUNC.  A regular expression over symbol names
      can match the generic C version, which the system may never run.  Our
      pipeline lifts the variant the machine selects (for example
      \code{\_\_strlen\_avx2}), found by address.  The skill could say how
      to choose.
\end{itemize}

\section{Discussion}\label{sec:discussion}

\paragraph{Complementary properties.}
OSS-Verify proves a weak property (no detectable error, agreement with the
plain build) about entire programs, mechanically, at large scale.  We prove
strong properties (functional correctness) about small units, with
specifications that need judgment.  The two connect naturally.  Inside
OSS-Verify proofs, agents already prove functional specifications of C
library routines such as musl's \code{memcpy}, because whole-program
reasoning needs them.  A verified, reusable library of such specifications
is what our project produces.

\paragraph{What Picinae could adopt.}
\begin{itemize}[nosep]
\item A checked release discipline: no \code{Admitted}, an axiom allowlist,
      and fingerprinted statements.
\item Statements generated from existing artifacts.  Test suites and fuzz
      harnesses could seed specifications that the model then strengthens.
\item An explicit fault semantics, so that a theorem can also rule out
      faults instead of being silent about them.
\end{itemize}

\paragraph{What the Picinae approach offers.}
\begin{itemize}[nosep]
\item Functional specifications, readable against the documentation.
\item An architecture-neutral IL: the same infrastructure already serves
      ARM, RISC-V and i386 in other Picinae examples.
\item Proofs small enough for a human to read and maintain.
\item Specification bugs found in a production C library (\code{sigisemptyset},
      \code{a64l}) by comparing the binary with its man page.  Crash-freedom
      statements cannot express such bugs.
\end{itemize}

\paragraph{Open questions.}
\begin{itemize}[nosep]
\item Can a model scale functional-correctness proofs the way OSS-Verify
      scaled agreement proofs, given that each specification needs review?
\item How much of the lifter belongs in the trusted base?  Validating
      lifted IL against an independent ISA semantics would close the gap
      with LPS's in-logic decoding.
\item Our slowest proofs are limited by kernel checking time, not by the
      agent.  What proof engineering would bring them to \code{Qed}?
\end{itemize}

\begin{thebibliography}{9}
\bibitem{anthropic-post}
Anthropic.
\newblock LLMs are making formal software verification feasible at a much
  larger scale.
\newblock \url{https://www.anthropic.com/research/formal-verification},
  October 2026.

\bibitem{oss-verify}
Theorem.
\newblock OSS Verify: theorem statements about OSS-Fuzz fuzz targets as
  compiled RISC-V binaries, and machine-checked Rocq proofs of them.
\newblock \url{https://github.com/theorem-labs/oss-verify-proofs}, 2026.

\bibitem{lps}
Theorem.
\newblock rocq-linux-process-semantics (LPS).
\newblock
  \url{https://github.com/theorem-labs/rocq-linux-process-semantics}.
\end{thebibliography}

\end{document}
